% First draft. Standalone companion to Machine Societies, not an extension of its theorem inventory.
\documentclass[11pt,letterpaper]{article}
\usepackage[margin=1in,headheight=15pt]{geometry}
\usepackage[T1]{fontenc}
\usepackage[utf8]{inputenc}
\usepackage{newtxtext}
\usepackage{microtype}
\usepackage{amsmath,amsthm,mathtools}
\usepackage{newtxmath}
\usepackage{booktabs,array,longtable,tabularx,graphicx,tocloft}
\usepackage{enumitem,fancyhdr,titlesec,needspace,placeins}
\usepackage{xcolor}
\definecolor{ink}{HTML}{172A44}
\definecolor{muted}{HTML}{5B6778}
\usepackage[pdfencoding=auto,psdextra]{hyperref}
\usepackage{bookmark}
\hypersetup{hidelinks,pdftitle={Machine Understanding: Representation, Explanation, and Transfer},pdfauthor={Komorebi P.},pdfsubject={First theoretical draft on reusable representations, explanatory transfer, and correction},pdfkeywords={machine understanding, representation, causal abstraction, explanation, transfer, revision, realization calculus}}
\pagestyle{fancy}\fancyhf{}
\fancyhead[L]{\small\textsc{Machine Understanding}}
\fancyhead[R]{\small First draft, 2026}
\fancyfoot[C]{\thepage}
\renewcommand{\headrulewidth}{0.3pt}
\titleformat{\section}{\large\bfseries}{\thesection}{0.65em}{}
\titleformat{\subsection}{\normalsize\bfseries}{\thesubsection}{0.65em}{}
\titlespacing*{\section}{0pt}{2.5ex plus .4ex}{1.2ex}
\titlespacing*{\subsection}{0pt}{1.8ex plus .3ex}{.8ex}
\setlength{\parindent}{1.2em}\setlength{\parskip}{.24em}
\setlength{\emergencystretch}{2em}
\clubpenalty=10000\widowpenalty=10000\displaywidowpenalty=10000
\setlength{\textfloatsep}{18pt plus 2pt minus 2pt}\raggedbottom
\setlist{leftmargin=1.5em,topsep=.4em,itemsep=.25em,parsep=0pt}
\setcounter{tocdepth}{1}\setlength{\cftbeforesecskip}{5pt}
\theoremstyle{plain}
\newtheorem{theorem}{Theorem}[section]
\newtheorem{proposition}[theorem]{Proposition}
\theoremstyle{definition}
\newtheorem{definition}[theorem]{Definition}
\newtheorem{counterexample}[theorem]{Counterexample}
\newtheorem{example}[theorem]{Example}
\newcommand{\X}{\mathcal X}\newcommand{\Z}{\mathcal Z}
\newcommand{\U}{\mathcal U}\newcommand{\Q}{\mathcal Q}
\newcommand{\Y}{\mathcal Y}\newcommand{\C}{\mathcal C}
\newcommand{\M}{\mathcal M}\newcommand{\E}{\mathbb E}
\newcommand{\Ffive}{\mathbb F_5}\newcommand{\TV}{d_{\mathrm{TV}}}
\newcommand{\img}[1]{\operatorname{im}(#1)}
\newcommand{\Eq}{\equiv_{\C}}\newcommand{\supp}{\operatorname{supp}}
\newcommand{\code}[1]{\texttt{\small #1}}
\newcolumntype{P}[1]{>{\raggedright\arraybackslash}p{#1}}
\begin{document}
\hypersetup{pageanchor=false}
\begin{titlepage}
\thispagestyle{empty}
\vspace*{.25in}
\noindent{\color{muted}\small\textsc{Operation Komorebi}\hfill\textsc{Theoretical research manuscript}}\par
\vspace{.65in}
\noindent{\color{ink}\fontsize{35}{40}\selectfont\bfseries Machine Understanding\par}
\vspace{.19in}
\noindent{\fontsize{17}{22}\selectfont Representation, Explanation,\\and Transfer\par}
\vspace{.28in}
\noindent{\color{ink}\rule{\textwidth}{.85pt}}\par
\vspace{.16in}
\noindent\textbf{Komorebi P.}\hfill\textit{First theoretical draft}\\
\noindent 8 October 2026\hfill Version 0.1.0\par
\vspace{.46in}
\noindent{\color{ink}\large\bfseries Abstract}\par\smallskip
\noindent A correct answer does not by itself identify a reusable method, a faithful account of the computation, or an explanation another reasoner can use. This paper develops a contract-relative account of computational understanding through the representations that support prediction, intervention, transfer, and correction. A finite controlled system provides the mathematical core. Behavioral equivalence characterizes the smallest exact reusable state, while an enlarged operation family exposes distinctions that an earlier representation could legitimately discard. We separate the availability of information from its causal use and require the cost and information access of encoders, decoders, and alignment maps to remain visible. Exact and approximate transport results describe when a method survives a change of representation. Explanations are treated as artifacts interpreted by a specified recipient, with explicit distinctions between subject understanding and faithfulness to a source computation. A correction theorem identifies when translating evidence preserves revision and when it only preserves a weaker set of candidate models. A three-register instrument connects these results in one worked example. The contribution is an integrated specification and a set of conditional arguments, not a universal test of understanding, an empirical claim about a deployed model, or a new account of subjective experience.
\vfill
\noindent\rule{\textwidth}{.4pt}\par\smallskip
\noindent{\color{muted}\small\textbf{Draft status.} Self-contained prose proofs and bounded exact companion checks. Established foundations and new exposition are distinguished. Independent review and formal proof-assistant verification remain open.}
\end{titlepage}
\hypersetup{pageanchor=true}
\pagenumbering{roman}
\begingroup\small\setlength{\parskip}{0pt}\tableofcontents\endgroup
\clearpage
\section*{Reading the draft}\addcontentsline{toc}{section}{Reading the draft}
The paper asks what a computational system can carry from one question to the next, from one representation to another, and from an accepted model to a corrected one. Its central object is a usable representation together with its operations and limits. It is not a score assigned to a model after one successful answer.

A conceptual reading can follow Sections~\ref{sec:intro}--\ref{sec:contract}, the interpretations after the formal results, and Section~\ref{sec:limits}. A mathematical reading proceeds through the finite quotient, transport, explanation, and correction results in order. Appendix~\ref{app:trace} gives one complete numerical trace. The source package includes exact reference programs, a statement index, a claim ledger, and a review record.

\paragraph{What is inherited and what is proposed.}
State minimization, causal abstraction, hypothesis testing, counting bounds, and error propagation have substantial established literatures. The proofs below make the particular assumptions used here inspectable. Their presence is not a claim of priority for those mathematical ingredients. The proposed integration is to require a claim about understanding to identify its target, its realized use, its admissible changes of representation, its recipient, its correction procedure, and its total resource boundary.

\paragraph{Relation to Operation Komorebi.}
\emph{Machine Societies} studies persistent plural institutions and the preservation of their operations and commitments~\cite{societies}. This companion changes the object of study. It concerns methods and representations that one computational reasoner can use, explain, transfer, and revise. It uses the same realization perspective without treating institutional continuity as evidence of individual understanding. No internal agent experiments are reported here.

\paragraph{Conventions.}
The main systems are finite, deterministic, and fully specified. A separate stochastic result states its own assumptions. A metric transport result also applies outside finite spaces, with explicit uniform bounds. A finite check supports its implemented construction, not the general theorem. Mathematical validity inside a model is distinct from evidence that an external system realizes that model.

\begin{center}\small
\begin{tabularx}{\textwidth}{@{}lX@{}}\toprule
Symbol & Meaning\\\midrule
$\C$ & Declared family of states, operations, and future questions.\\
$F_u$, $F_w$ & One implemented operation, or operations in the order specified by a word $w$.\\
$r:\X\to\Z$ & Retained representation. Its acquisition procedure is a separate obligation.\\
$\bar F_u$, $\bar q$ & An update and a query decoder operating on the retained state.\\
$T$, $\tau$ & Translation of representations and the corresponding operation names.\\
$V\subseteq\M$ & Current set of candidate models, not a probability distribution.\\
$C_{e,o}(V)$ & Candidates retained after experiment $e$ returns observation $o$.\\
$\phi$, $\eta$ & Translation of candidate models and translation of observations.\\
$\TV(P,Q)$ & $\tfrac12\sum_x|P(x)-Q(x)|$ on a finite outcome space.\\\bottomrule
\end{tabularx}\end{center}
\clearpage\pagenumbering{arabic}

\section{What a correct answer leaves open}\label{sec:intro}
A reasoner is asked what a device will display after two operations. It gives the right answer. That success is a fact worth preserving. It does not yet tell us whether the reasoner has a method that will work after three operations, whether the order of the operations matters, or whether the same method could be expressed in another notation. It also does not tell us whether the explanation offered afterward describes the computation that produced the answer.

These are not objections to correctness. They are different questions about what correctness rests on and what it makes possible. A proof can establish a conclusion from premises without teaching a recipient how to select the relevant premises in a new case. A useful model can guide new predictions without being a faithful description of the particular implementation that supplied it. A fluent explanation can be understandable while failing to identify the mechanism it claims to explain. The distinction between persuasive and faithful explanations is explicit in interpretability research~\cite{jacovi}.

This paper investigates a narrower and more tractable question than whether a machine understands in every human sense of the word. It asks when a computational realization possesses a representation that remains useful across a declared family of operations, questions, translations, and corrections. We call this \emph{reusable computational understanding} as an operational research proposal. It is relative to a domain and a contract. It is not a metaphysical property inferred from a single formal condition.

The central thesis is that a representation earns its explanatory role through what a system can do with it. It must retain distinctions needed for future questions, support the operations it claims to summarize, and expose the conditions under which its use should change. When it is communicated, a recipient must be able to recover something useful from it. When it is attributed to a particular mechanism, the attribution requires evidence about that mechanism. A representation that meets one of these obligations need not meet the others.

\subsection{Three separations organize the argument}
The first separation is between an answer and a reusable state. The answer to today's question can discard distinctions that tomorrow's permitted operation makes important. Understanding in the present sense therefore concerns a family of possible continuations, not only a completed input and output pair.

The second is between a correct account of a subject and a faithful account of a solver. An algebraic derivation can explain a numerical result even when the original solver used a table. It would be wrong to infer that the solver executed the derivation merely because the derivation is valid. Subject understanding and process faithfulness require different witnesses.

The third is between successful transport and preserved correction. Two languages may translate every currently expressible observation while one language merges distinctions needed to reject a mistaken model. Agreement in the translated language is then weaker than agreement about what has been learned. The correction theorem in Section~\ref{sec:correctiontransport} makes that loss exact.

\subsection{Position among existing theories}
Several neighboring theories already supply essential pieces. Computational mechanics studies minimal predictive state representations~\cite{shalizi}. Controlled state aggregation studies equivalences that preserve transition structure~\cite{givan}. Information bottleneck methods formalize compression relative to a specified relevance variable~\cite{tishby}. These are not interchangeable claims, but they establish that useful compression must be assessed against what it promises to preserve.

Causal abstraction gives a stronger account of agreement under interventions. Beckers and Halpern distinguish increasingly restrictive abstraction conditions~\cite{beckers}. Geiger and colleagues develop this perspective as a foundation for mechanistic interpretability, including transformations of mechanisms~\cite{geiger}. Our finite commuting equations occupy a deliberately restricted part of that subject. They do not replace a general theory of causal models.

There are also reasons to constrain the maps used in an interpretation. Hewitt and Liang show why a trained probe's success must be assessed against the work the probe itself can perform~\cite{hewitt}. Sutter and colleagues examine how sufficiently expressive alignment maps can make some causal-abstraction claims uninformative~\cite{sutter}. The constructive theory and this criticism address different aspects of the problem. We retain both lessons: intervention compatibility matters, and the translator must not be allowed to hide the entire explanation inside an unrestricted computation.

Explanation has a recipient as well as a referent. Miller's synthesis emphasizes the contextual and social features of explanation~\cite{miller}. Gentner's structure-mapping account places relational organization, rather than surface resemblance alone, at the center of analogy~\cite{gentner}. Here we formalize only a limited recipient-relative property: an artifact enables specified future operations in a recipient with stated background resources. That property does not exhaust the reasons a human explanation can succeed or fail.

Finally, learning through informative queries and counterexamples has a substantial theoretical history, including Angluin's work on regular languages~\cite{angluin}. Our correction model is much simpler. It is a finite filter on candidate models, chosen so that the interaction between evidence and translation can be proved directly.

\subsection{The intended contribution}
The paper does not propose that every useful representation is new, that a symbolic method is inherently superior to a learned one, or that a large table can never support understanding. It proposes a common account of the obligations that such representations must satisfy. The exact finite results locate information requirements and compatibility conditions. The counterexamples show why isolated successes cannot silently substitute for a combined claim.

This is also a limit on the theory itself. A finite contract can be satisfied by many implementations, including tables. To distinguish them, one must add relevant constraints on acquisition, causal organization, storage, computation, or transfer. Those constraints cannot be recovered from a definition after the results are known.

\section{A small instrument with a hidden register}\label{sec:instrument}
Consider a mathematical instrument with three registers
\[
x=(a,b,c)\in\Ffive^3.
\]
Each register contains one of $0,1,2,3,4$. All arithmetic in this example is modulo five. The display reports
\[
q(a,b,c)=a+b.
\]
There are $125$ possible complete states, but only five possible readings. The registers and their operations are fully specified analytical objects. They are not measurements from an agent experiment.

Initially, two operations are available:
\begin{equation}\label{eq:instrument-ops}
H(a,b,c)=(a+b,b,c),\qquad W(a,b,c)=(b,a,c).
\end{equation}
The first adds the second register into the first. The second exchanges the first two registers. The third register is preserved by both and does not affect the display. The operation order is public. A prediction concerns what the display would show after a named finite sequence of these operations.

\subsection{Knowing the current display is insufficient}
A record containing only $s=a+b$ answers the current question exactly. It cannot, in general, predict the display after $H$. The two complete states
\[
x=(0,0,0),\qquad y=(1,4,0)
\]
both display zero. After $H$, their displays are zero and four. Any predictor receiving only the common old display must return the same answer for both, so it cannot be exact in both cases.

A sufficient retained state for the original operation family is
\[
z=(s,b)=(a+b,b).
\]
Its updates are
\begin{equation}\label{eq:display-coordinates}
\bar H(s,b)=(s+b,b),\qquad \bar W(s,b)=(s,s-b),
\end{equation}
and its displayed output is its first coordinate. These equations retain the information needed for every finite $H/W$ continuation without retaining $c$.

The representation is also minimal in its number of states. The present reading $s$ and the reading $s_H$ after a hypothetical $H$ determine
\[
b=s_H-s,\qquad a=2s-s_H.
\]
Two different pairs $(a,b)$ therefore differ on at least one of these questions. Every exact reusable representation for this contract needs at least $25$ distinguishable values. Section~\ref{sec:quotients} gives the general argument.

\subsection{A new operation changes what must be known}
Now admit
\begin{equation}\label{eq:new-operation}
J(a,b,c)=(a+c,b,c).
\end{equation}
The old representation remains correct for $H$ and $W$. It is no longer sufficient for the enlarged contract. States with the same $(a,b)$ but different $c$ now produce different readings after $J$.

The additional reading $s_J$ yields
\[
c=s_J-s.
\]
Together with the earlier distinctions for $a$ and $b$, the enlarged question family distinguishes all $125$ complete states. The useful state counts are therefore five for the current display, $25$ for reusable $H/W$ prediction, and $125$ after $J$ is admitted.

\begin{figure}[htbp]
\centering\includegraphics[width=.98\linewidth]{figures/representation_scope.pdf}
\caption{The same instrument supports different minimal representations under different contracts. The numbers count distinguishable states, not scores of intelligence. A new permitted operation makes a previously irrelevant register relevant.}\label{fig:scope}
\end{figure}

This is not a case of an old theorem becoming false. The old theorem had a smaller operation family. What failed was an unlicensed extension of its conclusion. The distinction matters whenever a learned method is reused under a change in task, environment, or intervention.

\subsection{Correction requires evidence, not a renamed variable}
A predictor that has retained only $(s,b)$ cannot recover $c$ merely by adding a third field to its data structure. Conditional on a fixed $(a,b)$ and a uniformly distributed $c$, every possible display after $J$ is equally likely. Any guess based only on the old record succeeds with probability $1/5$. Randomizing the guess does not improve that probability.

An actual execution of $J$, followed by a new reading, supplies the missing distinction. Because $J$ changes the instrument, the recovered state is the state \emph{after} that operation:
\begin{equation}\label{eq:recover-after-J}
(a',b',c')=(s_J-b,\ b,\ s_J-s).
\end{equation}
The old state can then be reconstructed because this particular operation is invertible. Neither reversibility nor a free observation should be assumed for other instruments.

The example separates four achievements: reading an answer, retaining a method, recognizing a contract change, and acquiring evidence that repairs the resulting information gap. No amount of confidence in the original method supplies the fourth.

\section{A contract for reusable understanding}\label{sec:contract}
Let $\X$ and $\Y$ be finite nonempty sets. Let $\U$ be a finite operation alphabet, possibly empty, and let $\Q$ be a finite nonempty family of queries $q:\X\to\Y$. Each $u\in\U$ has an implemented total transition $F_u:\X\to\X$. The deterministic mathematical contract is
\[
\C=(\X,\U,(F_u)_{u\in\U},\Q,\Y).
\]
For a word $w=u_1\cdots u_n$, define $F_w=F_{u_n}\circ\cdots\circ F_{u_1}$. The empty word acts as the identity. Thus $q(F_w(x))$ asks a question after a stated continuation.

Totality is a simplifying assumption. If an operation is sometimes unavailable, its admissibility must be included in the observable contract or the model must use partial transitions with explicit domain preservation. An unmodeled refusal cannot be treated as the same event as a permitted update.

\begin{definition}[Exact reusable representation]\label{def:representation}
An exact reusable representation for $\C$ is a surjection $r:\X\to\Z$, updates $\bar F_u:\Z\to\Z$, and decoders $\bar q:\Z\to\Y$ such that, for all $x,u,q$,
\begin{equation}\label{eq:representation-contract}
r(F_u(x))=\bar F_u(r(x)),\qquad q(x)=\bar q(r(x)).
\end{equation}
\end{definition}
Induction immediately gives $q(F_w(x))=\bar q(\bar F_w(r(x)))$ for every finite word. The same initial record supports many future questions. It is not allowed to inspect the selected future query and replace itself with that query's answer before the challenge begins.

\subsection{A representation of the problem is not yet a representation in the solver}
The map $r$ specifies what information would suffice. It does not establish that a particular agent can acquire $r(x)$ from its sensors or that the agent uses this state. Those claims require an acquisition procedure, a realized computational path, and admissible interventions on that path.

For the instrument, an omniscient mathematical encoder can read $(a,b,c)$. A device restricted to the display cannot do so. It can obtain $b$ through a controlled $H$ observation, but that procedure costs an operation and changes the state. Any account of its understanding must include those effects rather than identify mathematical existence with operational access.

A practical contract therefore adds an information boundary, a family of recipient or decoder algorithms, resource limits, an intervention catalogue, and a revision policy. The finite theorems below solve particular obligations within that larger contract. They do not fill in unspecified sensors, access rights, or learning procedures.

\subsection{A profile rather than a single threshold}
We propose reporting a profile with separate entries for predictive adequacy, implemented use, intervention compatibility, transfer, recipient reconstruction, correction, and resource cost. Unknown entries remain unknown. A high prediction score must not silently fill the mechanism or correction entries.

The profile can describe a traditional program, a learned model, a hybrid system, or a collective procedure. It does not require language-like internal states. Nor does it establish an ordering of all systems from least to most understanding. Two systems can preserve different families of questions and be incomparable under the declared criteria.

\begin{counterexample}[A separate answer for each question is not one reusable record]\label{ex:joint}
Let $x=(a,b)\in\{0,1\}^2$ and let the two queries return $a$ and $b$. If the encoder knows which query will be asked, one bit suffices: retain that answer. If one record must be fixed before the query is selected and must answer both exactly, it needs four distinct values. A one-bit record has only two values and must identify two different answer profiles.
\end{counterexample}
The order of quantifiers is decisive. A different adequate representation for every question is not one adequate representation for their union. This is why a claimed method should be assessed with its encoder, updater, and decoder fixed across the family it promises to support.

\section{Behavioral compression and the price of new questions}\label{sec:quotients}
Define
\begin{equation}\label{eq:equivalence}
x\Eq y\quad\Longleftrightarrow\quad q(F_w(x))=q(F_w(y))\ \text{for all }q\in\Q,\ w\in\U^*.
\end{equation}
This is equality of all declared future answer profiles. It is an equivalence relation. It describes behavioral sufficiency under the contract, not equality of hidden mechanisms or identity outside that contract.

\begin{theorem}[The minimal exact reusable state]\label{thm:quotient}
The quotient $\X/{\Eq}$ carries an exact reusable representation for $\C$. Every other exact reusable representation $r:\X\to\Z$ refines this quotient: if $r(x)=r(y)$, then $x\Eq y$. There is a unique map $h:\Z\to\X/{\Eq}$ with $h(r(x))=[x]$. In particular,
\[
|\Z|\geq|\X/{\Eq}|.
\]
A fixed-length binary state needs at least $\lceil\log_2|\X/{\Eq}|\rceil$ bits. The quotient can be computed by finite partition refinement.
\end{theorem}
\begin{proof}
If $x\Eq y$, the empty word gives $q(x)=q(y)$ for every query. Moreover, $F_u(x)\Eq F_u(y)$ because any continuation after $u$ is a continuation in the original contract. Thus $[x]\mapsto[F_u(x)]$ and $[x]\mapsto q(x)$ are well-defined updates and decoders.

For any exact representation, the equations in~\eqref{eq:representation-contract} imply equality of all future decoded answers whenever $r(x)=r(y)$. This proves refinement. Surjectivity of $r$ gives existence and uniqueness of $h$ and the cardinality bound. The binary bound follows by counting available codewords.

For computation, let $P_0$ group states with identical current query vectors. Given a partition $P_k$, refine each block by the block of $F_u(x)$ for every $u$. Induction shows that $P_k$ identifies exactly the states agreeing on words of length at most $k$. Every strict refinement increases the block count. At most $|\X|-1$ strict refinements are possible. At a fixed point the partition is closed under all operations, hence under every finite word, and is exactly~\eqref{eq:equivalence}.
\end{proof}

This is a deterministic state-minimization result, presented in the language of reusable questions. It belongs to the established state-aggregation and predictive-state lineage~\cite{givan,shalizi}. Its significance here is not the discovery of quotients. It is the explicit identification of which future distinctions an asserted understanding is required to preserve.

\subsection{Minimal information is not minimal effort}
The theorem minimizes the number of distinct retained states. It does not minimize the work required to acquire the state, compute an update, verify a query, or discover the quotient. A smaller table can have harder entries to compute. A larger representation can make a frequently used transformation local and cheap.

Even uniqueness of the behavioral quotient does not supply a uniquely intelligible notation. The coordinates $(a,b)$ and $(s,b)$ describe the same $25$ behavioral classes in the original instrument. One makes the physical register update obvious. The other makes the readout immediate. Their difference is computational and explanatory rather than a difference in information content.

\subsection{The cost of expanding the contract}
Let $\C'$ have the same complete state space and contain every operation and query of $\C$, while possibly adding more. Write $\equiv_{\C'}$ for its finer behavioral equivalence. Take any representation $r$ adequate for $\C$. For each $z$, let $k_z$ be the number of $\equiv_{\C'}$ classes that meet the fibre $r^{-1}(z)$.

\begin{proposition}[Contract expansion and supplementary information]\label{prop:expansion}
The relation $\equiv_{\C'}$ refines $\equiv_{\C}$. Any supplementary code $s:\X\to\mathcal S$ for which $(r,s)$ supports exact reusable behavior under $\C'$ must satisfy
\[
|\mathcal S|\geq\max_z k_z.
\]
If every $\equiv_{\C'}$ class lies within a single fibre of $r$, this bound is attained by an exact reusable pair $(r,s)$. This condition holds when $r$ is the old minimal behavioral quotient. A supplementary value computed solely from $r(x)$ cannot repair a fibre meeting more than one expanded class.
\end{proposition}
\begin{proof}
Agreement on the larger continuation family implies agreement on the smaller one. Within a fixed fibre of $r$, an adequate pair $(r,s)$ must separate the $k_z$ new classes, giving the lower bound. For attainment under the additional condition, assign distinct labels $1,\ldots,k_z$ to the expanded classes within each old fibre and let $s(x)$ be its class label. The pair $(r(x),s(x))$ then identifies exactly the expanded behavioral class. The quotient update and query decoder supply its exact dynamics. Finally, any value computed only from $r(x)$ is constant on that fibre and separates none of its states.
\end{proof}

The attainment condition matters. An old representation can retain irrelevant extra coordinates that cease to be dynamically closed under new operations. For example, start with two bits $(a,b)$, no operations, a constant-zero query, and the unnecessarily detailed record $r(a,b)=a$. Add the swap operation while keeping the query constant. There is still only one expanded behavioral class, so $k_z=1$ in each old fibre. Yet updating the retained coordinate $a$ after the swap requires $b$. The information lower bound for answers is not sufficient to preserve arbitrary redundant coordinates. The minimal old quotient has no such problem.

In the instrument, each old $(a,b)$ class contains five expanded classes. An additional five-valued symbol is sufficient and necessary. A separately stored fixed-length symbol therefore needs three bits. This should not be confused with subtracting the rounded bounds for a joint $25$-state and $125$-state encoding, which are five and seven bits respectively. Separate buffers and joint recoding are different coding contracts.

\begin{counterexample}[A finite continuation test leaves an unbounded alternative]\label{ex:horizon}
For any fixed $h\geq0$, a one-operation machine can return zero through the first $h$ steps and return one at step $h+1$. A constant-zero machine agrees with it on all those earlier observations. Without a state bound or another restriction on the hypothesis family, success through horizon $h$ cannot establish agreement at every future horizon.
\end{counterexample}
This does not invalidate finite certification of a supplied finite model. It distinguishes exhaustive checking of a known transition system from identifying an unrestricted unknown system by a finite record.

\section{Available information is not implemented understanding}\label{sec:use}
The finite quotient answers an information question. A claim about a particular computational reasoner must answer an additional question: where is the useful representation realized, and how does it affect the computation?

A representation can contain a recoverable fact without the system using that fact. Conversely, a system can use information through a distributed or redundant encoding that a chosen simple probe fails to recover. The comparison of information structures in statistical decision theory provides a useful background for asking what a record could support~\cite{blackwell}. That question is different from identifying the route by which one implemented decision was made.

\subsection{A carrier can be readable and irrelevant}
Consider a solver state $m=(a,b,c,z)$, with $z=a+b$. Suppose its output circuit reads $a$ and $b$ directly and ignores $z$. A decoder applied to $z$ predicts the output perfectly. Replacing $z$ while keeping the actual circuit inputs unchanged has no effect. The code $z$ is available to an observer but is not the operative route used by this solver.

Now alter the implementation so that the output reads $z$, while a repair circuit recomputes $z$ from $a,b$ before every output. An attempted replacement may again have no observable effect, but for a different reason: the affected carrier is restored. A null intervention cannot identify which of these implementations is present unless its timing and the remaining information routes are controlled.

These examples justify neither a universal ablation test nor a demand that every useful representation be uniquely necessary. They show why an intervention must specify what changes, what remains accessible, and what the hypothesis predicts afterward. A proof about the instrument's state does not substitute for an intervention on the solver's implementation.

\begin{proposition}[The identification limit of a transcript]\label{prop:identification}
Suppose two proposed implementations induce complete transcript laws $P_0$ and $P_1$ on a finite transcript alphabet under a specified permitted experiment, including any adaptive choices made by its evaluator. With equal prior probability, the greatest identification accuracy available from that transcript is
\[
\frac12\bigl(1+\TV(P_0,P_1)\bigr).
\]
In particular, identical transcript laws permit no accuracy above $1/2$.
\end{proposition}
\begin{proof}
Let $d(o)\in[0,1]$ be the probability of guessing implementation 0 after transcript $o$. Accuracy is
\[
\frac12\sum_o\bigl(P_0(o)d(o)+P_1(o)(1-d(o))\bigr)
=\frac12+\frac12\sum_o(P_0(o)-P_1(o))d(o).
\]
The sum is maximized by setting $d(o)=1$ on positive differences and zero on negative ones. The positive differences sum to total variation. This choice attains the bound.
\end{proof}
The experiment may include interventions. The proposition says only that the \emph{permitted complete transcript} bounds identification. Adding a new sensor or an admissible internal intervention changes that transcript law and can change the limit.

\subsection{The interpreter must not secretly do all the work}
Take any bijection $p:\X\to\Z$ of a finite state space. Define
\[
\bar F_u=p\circ F_u\circ p^{-1},\qquad \bar q=q\circ p^{-1}.
\]
All exact representation equations hold. The labels in $\Z$ may nevertheless be arbitrary, and the translated updates may require a complete lookup table. Exact representability alone does not explain why the representation is easier to acquire, use, or communicate.

This does not make the construction false. It makes the claim narrow. A table can be a valid finite model. Its interpretive value depends on the allowed questions and costs. The mistake would be to infer that any such recoding identifies a compact or causally privileged concept inside a solver.

We therefore require an \emph{interface account}. The encoder, decoder, intervention implementation, and translator are part of the evaluated system. Their information access is declared, their computation is charged, and any training or target-specific background they receive is recorded. A decoder that is free to solve the whole task is not evidence that the state it reads supplied the method. Control tasks for probes and the critique of unrestricted alignment maps are especially relevant here~\cite{hewitt,sutter}.

Restrictions can concern circuit size, memory, communication rounds, a finite program language, or an explicit family of allowed maps. There is no universally correct restriction. A linear map is not automatically the right representation of a nonlinear concept, and a nonlinear map is not automatically uninformative. The appropriate restriction must be selected with the scientific question, not retrofitted to distinguish preferred results.

\subsection{A bounded stochastic extension}
The distinction between a current answer and a reusable state also survives stochastic dynamics. Let $K_u(x,x')$ be finite transition kernels and let $r:\X\to\Z$ be surjective. Assume the current query values are constant on fibres of $r$.

\begin{proposition}[Uniform controlled stochastic reuse]\label{prop:lumpability}
A single reduced kernel $\bar K_u$ reproduces every projected one-step law from every complete state if and only if, whenever $r(x)=r(y)$,
\[
\sum_{v:r(v)=z'}K_u(x,v)=\sum_{v:r(v)=z'}K_u(y,v)
\quad\text{for every }u,z'.
\]
With matched initial projected laws, the reduced model preserves the law of every finite sequence of reduced states and queries under any policy depending only on the reduced history.
\end{proposition}
\begin{proof}
A common projected row is necessary by considering point-mass initial states. If the rows agree, use their shared value to define $\bar K_u(r(x),z')$. Conditional on any reduced history, the complete state has a distribution supported on the current fibre. Averaging equal projected rows gives the same reduced transition law. Induction yields equality of the finite history laws, including the deterministic query values. A common randomized history policy preserves the same argument after conditioning on its selected action.
\end{proof}
This is the controlled lumpability condition underlying stochastic state aggregation~\cite{givan}. The policy qualification is essential. A controller that secretly observes a discarded variable can choose different actions in states that the reduced description identifies. Small prediction error on one observed trajectory distribution is not a substitute for the uniform condition.

\section{Transporting a method rather than matching an answer}\label{sec:transport}
An analogy is useful when relations and operations survive the comparison, not simply when two descriptions use similar words. Structure-mapping theory makes relational correspondence central to analogy~\cite{gentner}. We study a more limited mathematical question: when can an already specified method be translated without changing the answers to its declared continuation questions?

Let a source representation have state space $Z$, updates $f_u$, and query decoders $q$. Let a target have state space $W$, updates $g_v$, and decoders $q'$. A map $\tau$ translates operation names. Query pairs $(q,q')$ and a representation map $T:Z\to W$ are fixed in advance. The operation names alone do not establish that the target implements the intended source operation.

\begin{theorem}[Exact transport of reusable operations]\label{thm:transport}
Suppose, for every source operation and matched query,
\begin{equation}\label{eq:transport-exact}
T f_u=g_{\tau(u)}T,\qquad q'T=q.
\end{equation}
Then, for every source state and every finite source word $w$,
\[
T f_w=g_{\tau(w)}T,\qquad q'(g_{\tau(w)}(Tz))=q(f_w(z)).
\]
Two such transports compose, with their operation and query correspondences composed as well. No injectivity assumption on $T$ is needed for this forward guarantee.
\end{theorem}
\begin{proof}
The empty word is immediate. If the identity holds for $w$, then
\[
T f_u f_w=g_{\tau(u)}T f_w
=g_{\tau(u)}g_{\tau(w)}T.
\]
Induction proves the word identity, and applying the matched decoder proves the answer identity. For a second transport $S$ with operation map $\sigma$, substitution gives $STf_u=h_{\sigma(\tau(u))}ST$ and the corresponding decoder identity. This is the same condition for the composite.
\end{proof}
This is an intertwining argument. Related composition questions for interventional causal models are studied in greater generality by Rischel and Weichwald~\cite{rischel}. Here the theorem guarantees forward answers only from states in the image of $T$ and for translated source operations. It does not guarantee coverage of all target states or all target operations. Surjectivity, reverse translation, and recovery of discarded source distinctions are separate requirements.

\subsection{The instrument in display coordinates}
For the original instrument quotient, let the source state be $(a,b)$ and set
\[
T(a,b)=(a+b,b).
\]
The target updates in~\eqref{eq:display-coordinates} satisfy~\eqref{eq:transport-exact}. The source decoder $a+b$ becomes the target's first coordinate. Thus a method can be carried between a register description and a display-centered description without relearning the meaning of every operation sequence.

\begin{figure}[htbp]
\centering\includegraphics[width=.96\linewidth]{figures/transport_square.pdf}
\caption{One verified operation square for the instrument. Translating and then applying the target update produces the same state as applying the source update and then translating. The equality concerns the declared $H$ operation, not arbitrary interventions.}\label{fig:transport}
\end{figure}

\begin{counterexample}[Correct current answers with a wrong transported operation]\label{ex:badtransport}
Retain the correct map $T(a,b)=(a+b,b)$ and the correct display decoder. Now incorrectly define the target swap by $\widetilde W(s,b)=(s,b)$. The current display and the display immediately after a swap are still correct. Starting from $(a,b)=(1,3)$, however, the source sequence $W$ then $H$ ends with display zero. The incorrect target sequence ends with display two. Matching a restricted collection of outputs did not preserve the method.
\end{counterexample}
The error becomes visible only when the altered internal state is used by a later operation. That is precisely why reusable understanding requires continuations rather than a list of isolated answers.

\subsection{Approximate transfer accumulates error}
Exact equalities are not always available. For this result let the target state space carry a metric $d$. Suppose each target update is $L$-Lipschitz for a common $L\geq0$, and suppose
\[
d(Tf_u(z),g_{\tau(u)}(Tz))\leq\epsilon
\]
uniformly on every state and operation that the contract permits. Let a target query be $M$-Lipschitz with real-valued output and satisfy $|q'(Tz)-q(z)|\leq\delta$.

\begin{proposition}[Finite-horizon transfer budget]\label{prop:error}
For matched source and target operation sequences, start from $z_0,w_0$ with $d(Tz_0,w_0)\leq e_0$. After $n$ operations,
\begin{align*}
d(Tz_n,w_n)&\leq L^n e_0+\epsilon\sum_{j=0}^{n-1}L^j,\\
|q(z_n)-q'(w_n)|&\leq\delta+M\left(L^n e_0+\epsilon\sum_{j=0}^{n-1}L^j\right).
\end{align*}
The empty sum is zero and $L^0=1$, including when $L=0$. At $L=1$ the state bound is $e_0+n\epsilon$.
\end{proposition}
\begin{proof}
Writing $e_n=d(Tz_n,w_n)$, insert $g_{\tau(u)}(Tz_n)$ and use the triangle inequality and the Lipschitz bound to obtain $e_{n+1}\leq\epsilon+Le_n$. Iterating gives the stated expression. Insert $q'(Tz_n)$ between the two outputs and use the decoder mismatch and its Lipschitz constant.
\end{proof}
The bound concerns the matched operation sequence. It does not automatically cover two controllers that select different operations after slightly different observations. Adaptive choices require a matching policy contract or a separate stability argument for the controller.

\begin{figure}[htbp]
\centering\includegraphics[width=.94\linewidth]{figures/error_accumulation.pdf}
\caption{Exact values of the recurrence $e_{n+1}=Le_n+1/100$, with $e_0=0$, illustrate the transfer bound. The curves are analytical constructions, not measured model errors. A small local defect can remain bounded, accumulate linearly, or be amplified.}\label{fig:error}
\end{figure}

A map can therefore be accurate for a short continuation and unsuitable for a longer one. The required tolerance belongs to the task, not to a universal declaration that the two representations are close. An approximate relation with one tolerance also need not be transitive. A chain of individually accepted translations may exceed the original error budget.

\subsection{Local representations form an atlas only with transition rules}
One representation may simplify prediction while another simplifies correction or proof. A useful system need not select one notation forever. It can maintain a family of representations with declared domains and translators. We call this an atlas only in the modest sense of a set of local descriptions connected by checked transition maps.

The domain guards matter. A method valid on reversible operations may fail after an irreversible reset. A chart switch may require evidence absent from the current chart. A common output name does not establish a common unit or query semantics. The system must know which transitions are supported and must retain the unresolved discrepancies when a transition is partial.

There is no general theorem here that more representations improve reasoning. Additional charts consume memory and create more opportunities for a mistaken translation. The research target is a family whose additional coverage or lower repeated cost justifies its construction and maintenance. Exact intertwining supplies one part of that analysis, not its economic conclusion.

\section{An explanation belongs to a recipient}\label{sec:explanation}
An explanation can be a correct description and still fail to help a particular recipient. It may presume unfamiliar notation, omit a necessary background result, or encode all the difficult work in a phrase the recipient cannot execute. Conversely, a deliberately simplified explanation can be useful for a limited question family without being a complete description of the source mechanism.

We therefore distinguish two contracts. A \emph{subject explanation} enables the recipient to reason correctly about specified aspects of the subject. A \emph{process explanation} additionally makes a faithful claim about how a particular solver reached its result. A process explanation must satisfy an implementation attribution as well as a subject-level or transcript-level account. Neither fluent presentation nor a valid derivation alone establishes the attribution~\cite{jacovi}.

\begin{definition}[Recipient-relative explanatory adequacy]\label{def:explanation}
Fix a recipient procedure $L$, background information $B$, resource allowance, and continuation family. An artifact $E$ is adequate for a target model $m$ when $L(B,E)$ yields a representation and operations that answer the required continuations for $m$ within the declared error and resource limits. The artifact's inputs, the recipient's prior access, and any target-specific background are included in the contract.
\end{definition}
This is an operational property of an artifact and a recipient. It does not equate instruction with human comprehension, and it does not require natural language. A diagram, proof, worked example, executable transformation, or combination can play the role of $E$.

\subsection{An exact communication lower bound}
Let $\Theta$ be a finite family of target models. Relative to the required continuation questions, divide these models into $N$ classes with distinct answer profiles. Use one fixed recipient background and one fixed deterministic recipient. The teacher can select a message but cannot change the recipient's background separately for each target without that change being counted as communication.

\begin{proposition}[The information an exact explanation must carry]\label{prop:explanation}
If a fixed-length $b$-bit artifact lets the recipient answer the declared questions exactly for every target in $\Theta$, then $2^b\geq N$. With target-dependent background $B_\theta$, the same argument gives
\[
2^b\geq\max_{B}\ N_B,
\]
where $N_B$ counts distinct answer profiles among targets sharing that exact background value.
\end{proposition}
\begin{proof}
Two different required answer profiles cannot produce the same background and message. The deterministic recipient would then make the same responses for both and fail on a distinguishing question. There are at most $2^b$ fixed-length messages. Apply this counting argument to the whole family, or separately within each background fibre.
\end{proof}
The bound is about information, not explanatory elegance. It does not establish that any shortest message is easy to understand or that the recipient can decode it efficiently. A compressed index into an enormous shared table can meet the bound while transferring little new method. The table's construction and availability must be charged to the background account.

\subsection{What the instrument explanation transfers}
For the original instrument, an adequate artifact can specify arithmetic modulo five, the representation $(s,b)$, the two update rules in~\eqref{eq:display-coordinates}, the output decoder, and the domain restriction excluding $J$. A recipient supplied with an initial $(s,b)$ can then answer every $H/W$ continuation by applying these operations.

That establishes a subject explanation for the specified recipient. It does not establish that a neural system that produced the artifact used $(s,b)$ internally. A separate process explanation might show that the system's actual intermediate states realize these coordinates and respond to corresponding interventions. These two certificates can support one another, but they answer different questions.

An especially useful explanation includes the failure boundary. The line ``the third register is irrelevant'' is too broad. The line ``the third register does not affect any display reachable through the current $H/W$ operation family'' is exact. The latter statement tells the recipient what should be reconsidered when a new operation arrives.

\subsection{Explanations can preserve methods without reproducing discovery}
A proof may give a shorter route than the route by which the result was found. A teacher may introduce coordinates that the discoverer never used. These are legitimate explanatory choices when their subject-level claims are correct. They become misleading only when presented as faithful records of the original reasoning without evidence.

This distinction avoids an unnecessary demand that every helpful explanation be an autobiography of a computation. It also prevents the opposite mistake of treating a helpful reconstruction as proof of the original mechanism. Human context and communicative purpose remain important beyond the formal recipient criterion~\cite{miller}.

\section{Learning where the current representation fails}\label{sec:revision}
A reusable representation should support correction without treating every new observation as a reason to discard everything. A valid method can have a restricted domain. A recorded premise can be false. A previously omitted variable can become relevant. These failures call for different revisions.

We distinguish three cases. In a \emph{state correction}, the operations remain valid but the current state estimate changes. In a \emph{model correction}, an assumed transition or observation law changes. In a \emph{contract expansion}, the requested operations or questions change. A representation can be entirely adequate before the third kind of change and inadequate afterward, as the instrument shows.

\subsection{A finite model of informative correction}
Let $\M$ be a finite nonempty set of deterministic candidate models. Let $\mathcal E$ be a finite set of permitted experiments. Each experiment $e$ has an outcome map $f_e:\M\to\mathcal O_e$. The current information is a candidate set $V\subseteq\M$. On observing $o$ from $e$, update it by
\begin{equation}\label{eq:filter}
C_{e,o}(V)=\{m\in V:f_e(m)=o\}.
\end{equation}
This is consistency filtering, not Bayesian inference and not a general model of scientific reasoning. An empty set is meaningful: the observation contradicts every currently retained candidate. It must not be replaced by an invented best answer.

An experiment is informative for $V$ when its outcome map is not constant on $V$. Models indistinguishable under every permitted experiment form an observational equivalence class. The aim can be to identify such a class rather than a uniquely named model.

\begin{proposition}[Finite correction under separating evidence]\label{prop:queries}
Suppose the true model $m_*$ belongs to the initial set $V_0$, every observation is exact, and the allowed experiment family separates any two distinct remaining observational classes. Repeatedly selecting an informative experiment preserves $m_*$ and reaches a single observational class after at most $|V_0|-1$ informative observations. If there are $N$ distinguishable target classes, every adaptive procedure identifying all of them using experiments with at most $b\geq2$ outcomes and at most $h$ observations must satisfy $b^h\geq N$.
\end{proposition}
\begin{proof}
The true model predicts each exact observation, so filtering never removes it. An informative experiment splits the current set into at least two nonempty outcome classes. Whatever true outcome occurs, at least one candidate is removed. The nonempty set can strictly shrink at most $|V_0|-1$ times. If two observational classes remain, the separating assumption supplies another informative experiment. For the second claim, an adaptive decision tree of depth at most $h$ and branching at most $b$ has at most $b^h$ leaves. Distinct required classes must reach different leaves.
\end{proof}
This elementary result shares the query-and-counterexample perspective of classical active learning~\cite{angluin}. It does not give an efficient method for discovering the candidate family, finding an inexpensive separating experiment, or dealing with inaccurate observations. Its upper bound is a count of informative observations, not a bound on wall-clock work or physical cost.

\subsection{A question can be more useful than another answer}
If all surviving candidates give the same answer to the next task, the system can act without identifying which candidate is true. If they disagree on a consequential operation, an experiment that distinguishes that disagreement may be more valuable than repeating familiar tasks on which all candidates agree.

The criterion is task-relative. Two models can differ on an inaccessible quantity and still support all permitted decisions. Conversely, a rarely asked question can distinguish models whose disagreement controls whether a future intervention is safe. The system should retain which distinctions have been settled and which have merely been irrelevant so far.

A request for evidence is also an operation with a realization. It may consume resources, change the system, or reveal only a biased part of the state. The finite filtering model presumes that the stated outcome maps already include those effects. For a changing instrument, one must use a controlled reset, a model of the full observation history, or a state update that tracks what the experiment changed. Independent repeated access to the same initial state is not free.

\subsection{Repairing the running example}
Before $J$ is admitted, all five possible values of $c$ are equivalent under the base contract. Treating them as one class is correct, not ignorance that the old task required the predictor to resolve. After $J$ is admitted, these classes split. The system can report that its old state no longer determines the answer.

The current record supplies $(s,b)$. Executing $J$ and observing $s_J$ determines $c=s_J-s$, while~\eqref{eq:recover-after-J} updates the affected first register. The corrected extended coordinates are $(s_J,b,c)$. Their operations are
\[
\bar H(s,b,c)=(s+b,b,c),\quad
\bar W(s,b,c)=(s,s-b,c),\quad
\bar J(s,b,c)=(s+c,b,c).
\]
The base method survives as the restriction of the enlarged method. Only the newly relevant distinction is acquired.

A different kind of failure would occur if the instrument secretly implemented $H(a,b,c)=(a+2b,b,c)$. That would contradict the old transition model rather than merely extend the operation catalogue. A correction record should say which assumption changed. Otherwise a sequence of patched answers can conceal a permanently incorrect account of the mechanism.

\subsection{Revision must preserve its own uncertainty}
The true-model assumption in Proposition~\ref{prop:queries} is a substantive premise. When it is not justified, the system must allow the possibility that every candidate is wrong. When observations are noisy, deleting a model after one conflicting record can remove the best explanation. A statistical version needs a likelihood model, calibrated tolerances, and an account of selection effects.

The finite framework nevertheless gives a useful discipline. A correction should identify the new evidence, the rejected candidates, the surviving claims, and the assumption linking the observation to the model. It should not rewrite an old uncertainty as though the corrected answer had always been established.

\section{When translation commutes with correction}\label{sec:correctiontransport}
A method can travel through a language that preserves its current answers while losing the evidence needed to improve it. This is a different obstruction from a wrong update equation. It concerns the interaction between two otherwise sensible operations: translating a model and filtering it against evidence.

Let $\M$ and $\mathcal N$ be finite model families, with model translation $\phi:\M\to\mathcal N$. For one source experiment $e$, let $f:\M\to\mathcal O$ give its outcomes. A corresponding target experiment has outcome map $g:\mathcal N\to\mathcal O'$. Let $\eta:\mathcal O\to\mathcal O'$ translate evidence. Assume the basic observation square commutes:
\begin{equation}\label{eq:evidence-square}
g(\phi(m))=\eta(f(m))\quad\text{for all }m\in\M.
\end{equation}
This condition says that a model's predicted evidence translates correctly. It permits $\eta$ to merge source outcomes.

Write $\phi[V]=\{\phi(m):m\in V\}$. Let $C_o(V)=\{m\in V:f(m)=o\}$ and $C'_{o'}(W)=\{n\in W:g(n)=o'\}$. Both are ordinary consistency filters.

\begin{theorem}[Correction transport and evidence reflection]\label{thm:correction}
Under~\eqref{eq:evidence-square}, for every $V\subseteq\M$ and every attainable source outcome $o\in f(\M)$,
\begin{equation}\label{eq:correction-inclusion}
\phi[C_o(V)]\subseteq C'_{\eta(o)}(\phi[V]).
\end{equation}
Equality holds for all such $V$ and $o$ if and only if $\eta$ is injective on the attainable outcome set $f(\M)$. The claim includes candidate sets contradicted by the observation, so that preservation of an empty update is part of the contract.
\end{theorem}
\begin{proof}
If $n\in\phi[C_o(V)]$, then $n=\phi(m)$ for some $m\in V$ with $f(m)=o$. Equation~\eqref{eq:evidence-square} gives $g(n)=\eta(o)$, proving inclusion.

If $\eta$ is injective on $f(\M)$ and $n\in C'_{\eta(o)}(\phi[V])$, choose $m\in V$ with $\phi(m)=n$. Then $\eta(f(m))=g(n)=\eta(o)$, so $f(m)=o$. Thus $n\in\phi[C_o(V)]$, proving equality.

Conversely, suppose two distinct attainable outcomes $a,b$ have $\eta(a)=\eta(b)$. Choose $m$ with $f(m)=b$, let $V=\{m\}$, and consider source observation $a$. The source update is empty, whereas the target update retains $\phi(m)$ by~\eqref{eq:evidence-square}. Uniform equality therefore fails.
\end{proof}
This is an exact set-theoretic result about direct images and filtering. The institutional or cognitive interpretation does not make it a new general theorem of set theory. Its value is to expose a condition that output translation alone omits: the translated evidence must preserve the distinctions by which a candidate is excluded.

The necessity statement deliberately tests observations that contradict the current candidate set. If those cases are excluded by a standing assumption that the current model is already correct, weaker conditions can suffice. That is a different and less demanding correction contract. A reasoner meant to discover that its model is wrong should not assume away the contradictory case.

\begin{counterexample}[A translated observation preserves both mistaken candidates]\label{ex:correction}
Take $\M=\{m_0,m_1\}$ with $f(m_0)=0$ and $f(m_1)=1$. Translate models bijectively to $\{n_0,n_1\}$, but translate both observations to a single symbol $*$. Let $g(n_0)=g(n_1)=*$. The observation square commutes. Starting from both candidates, source observation $1$ leaves only $m_1$, whose image is $\{n_1\}$. Translating first and then filtering by $*$ leaves both $n_0$ and $n_1$.
\end{counterexample}

\begin{figure}[htbp]
\centering\includegraphics[width=.98\linewidth]{figures/correction_square.pdf}
\caption{Correct translation of each predicted observation need not preserve correction. The lower evidence channel merges zero and one, so the target cannot remove the candidate excluded at the source. The theorem specifies when the inclusion becomes equality.}\label{fig:correction}
\end{figure}

The failure is not an inaccurate target answer in its coarse language. Both models really do predict $*$ there. The failure is the loss of a previously meaningful ability to discriminate. This is why a representation that can communicate conclusions may still be an inadequate carrier of scientific revision.

\subsection{The instrument explains the boundary}
For base experiments built only from $H$, $W$, and the display, the translation from $(a,b,c)$ to $(a,b)$ can preserve evidence without collapsing outcome symbols. Different values of $c$ produce the same permitted observations, so their identification loses no base correction power.

For $J$, that translation cannot preserve the full outcome map at all: two states with identical $(a,b)$ can have different readings after $J$. One could force a commuting square by translating all new readings into one coarse symbol. That would preserve the trivial coarse prediction, but it would destroy the ability to discover $c$. Enlarging the representation is a different response from weakening the question until the old representation appears adequate.

This distinction is important for evaluation. A system should not be rewarded for transfer by silently replacing a difficult target question with an easier translated one. The source and target question contracts, including evidence reflection when correction matters, must be fixed before the comparison.

\subsection{Composing corrective explanations}
Suppose an explanatory artifact carries a recipient's candidate set, translated experiments, and evidence interpretation. If each translation preserves the observation square and reflects the relevant outcomes, successive filters and translations commute by repeated application of Theorem~\ref{thm:correction}. The recipient can then discard the same translated candidates as the source after the specified observation sequence.

This conclusion is only about filtering the represented model families. It does not transfer an ability to invent a missing model, establish that the observations are trustworthy, or guarantee efficient implementation. An explanation can teach how to correct within a family while leaving the enlargement of that family as a separate problem.

\section{One witness for prediction, transfer, and revision}\label{sec:witness}
The separate results now have to refer to one realized system. A good state representation, an elegant translation, and a useful explanation do not establish a combined capability when they require different hidden information, different decoders, or incompatible resource assumptions.

\begin{definition}[Integrated adequacy record]\label{def:witness}
An integrated adequacy record for a declared task family identifies a common computational realization, its information boundary, a reusable representation and its updates, permitted intervention implementations, each claimed transport and recipient, the model and evidence maps used for correction, and the cost and scope of those operations. Every claimed guarantee refers to these same objects or to explicitly related replacements.
\end{definition}
This is a specification proposal, not a new primitive of cognition. Its purpose is to make the quantifiers visible. The record cannot use a powerful target-specific decoder to establish prediction, then a different decoder to establish transfer, while reporting the result as one reusable method unless the selection procedure is itself part of the method.

\subsection{What can be checked in the finite core}
For explicitly supplied finite tables, the equations in Definition~\ref{def:representation} can be checked on every state, operation, and query. Theorem~\ref{thm:quotient} then proves the guarantee for arbitrary finite continuations. Exact transfer requires corresponding one-step and decoder checks. Correction transport requires the observation square and its stated reflection condition. The recipient must actually construct and execute the declared procedure within its allowed information boundary.

Some obligations remain different in kind. The finite model can list which intervention is admissible, but relating that declaration to a physical device requires evidence about the device. A table can name a sensor reading without proving calibration. A certificate can bind a program version without establishing the semantic meaning of its inputs. The specification should preserve these distinctions rather than present one successful check as universal certification.

\begin{center}\small
\begin{tabularx}{\textwidth}{@{}P{1.35in}XX@{}}\toprule
Claim & Required object & What it does not supply\\\midrule
Reusable prediction & One retained state, updates, decoders, and a fixed future family. & Access to that state or evidence of its actual use.\\
Mechanistic attribution & A realized mapping and discriminating permitted interventions. & Every untested mechanism or unrestricted causal access.\\
Method transfer & Matched states, operations, queries, and coverage. & Reverse reconstruction or new target operations.\\
Recipient explanation & Fixed background, an artifact, an interpreter, and a budget. & Faithfulness to the source solver's actual computation.\\
Corrective transfer & Candidate and evidence maps that preserve filtering. & Model invention, noise calibration, or truth outside the family.\\
Resource feasibility & One accounting boundary for construction and repeated use. & A universal efficiency advantage.\\\bottomrule
\end{tabularx}\end{center}

\subsection{The complete instrument witness}
For the original contract, the instrument witness contains $r(a,b,c)=(a+b,b)$, updates~\eqref{eq:display-coordinates}, the first-coordinate decoder, and a statement that $c$ is outside the relevant future distinctions. It also names how an observer acquires the initial record. A recipient with that record and the equations can reproduce all declared predictions.

After $J$ is admitted, the witness changes. The old representation does not predict the new reading. An explicit $J$ observation supplies $c$ and updates the first coordinate. The enlarged state has the exact updates in Section~\ref{sec:revision}. The original $H/W$ explanations remain valid on their unchanged domain. New predictions require the acquired evidence, and any translated explanation must preserve that evidence's discriminating meaning.

These are connected properties of one analytical construction. They do not establish that an unspecified learning algorithm will discover the representation, infer the correct transition law, or choose the diagnostic operation. The point is to make those remaining research obligations identifiable rather than confuse them with the solved representation problem.

\subsection{A cost account for reusable methods}
For a sequential workflow with disjoint accounting, total cost includes acquiring information, discovering the representation, encoding state, performing updates, verifying claimed relationships, translating when needed, and interpreting the result. Parallel execution requires work and elapsed time to be recorded separately. Costs in different units should not be added without a stated conversion.

A simple amortization comparison illustrates the issue. Suppose a representation costs $D$ to construct, each later use costs $c$, and the matched baseline costs $b>c$ per use. The representation has lower total cost after $N$ uses exactly when
\[
D+Nc<Nb,\quad\text{or equivalently}\quad N>\frac{D}{b-c}.
\]
This accounting does not prove that the representation can be discovered for cost $D$, that the uses are comparable, or that all target instances share it. It locates where such evidence would enter.

Understanding need not always make computation faster. A more costly representation can make a proof inspectable, reveal a missing assumption, or support an intervention unavailable to the cheaper one. The relevant benefit must be stated. The theory should not smuggle efficiency into the meaning of understanding, nor infer efficiency from the existence of a short description.

\section{Limits, objections, and the next mathematical questions}\label{sec:limits}
The framework is intentionally operational, but that does not make it philosophically complete. A contract specifies selected consequences. It does not establish that those are all the consequences that matter. A formal model of a subject can be wrong, and a model of a reasoner's computation can be incomplete. A successful finite witness does not answer every question about meaning, intentionality, or experience.

\subsection{Does this define understanding as success?}
Not as a single success rate. The framework separates what a state can support, what a solver actually uses, what can be transported, what a recipient can reconstruct, and what survives correction. The counterexamples show that these properties differ. Nevertheless, every operational criterion remains a criterion for a declared set of abilities. The draft does not offer an independent essence of understanding beyond all such criteria.

The resulting restraint is useful. A system may possess strong reusable knowledge in one domain and weak correction behavior in another. A receiver may understand an explanation under a supplied notation while failing to translate it. A process may be correctly modeled by an outside analyst without being able to explain itself. These are distinct descriptions rather than contradictory verdicts about an indivisible trait.

\subsection{Can a table satisfy the theory?}
Yes, for some finite contracts. A table can encode a transition system and support exact future queries. It should not be declared incapable by terminology alone. The stronger questions concern its origin, size, use, reuse across a family, and behavior when the contract changes.

The recoding construction in Section~\ref{sec:use} makes this unavoidable. To distinguish a compact transferable method from a complete enumeration, one must specify the comparison class and account for the encoder, decoder, and shared background. Restricting those components is part of the research design. An assertion that a representation is meaningful is not a substitute for doing so.

\subsection{Do valid proofs establish understanding?}
A proof establishes the conclusion licensed by its premises and rules. The recipient may still lack a usable model of why those premises were appropriate or how the proof changes under a new problem. Conversely, a reasoner can use a correct model while being unable to produce a formal proof in a particular proof language.

Formal verification is therefore valuable but not interchangeable with the entire profile. The present proofs establish finite and metric statements under their assumptions. They do not certify that a deployed neural model realizes the maps, that a human reader has understood them, or that the task specification captures all relevant external facts.

\subsection{Finite models and a changing world}
The exact quotient assumes a fixed state space and a fixed operation family. The correction theorem assumes specified model families and outcome maps. The instrument's new operation enlarges one of those families in a controlled way. Real inquiry can also introduce new variables, new causal hypotheses, new measurement channels, and new forms of explanation.

There is no guarantee that the minimal state for every future extension is small. Preserving all conceivable future distinctions can collapse the compression objective back toward retaining the entire accessible state, and even that state may omit external information. A useful theory must allow partial coverage and a principled request for new evidence rather than promise unlimited generality.

\subsection{Four questions for the next revision}
\paragraph{Bounded nonvacuity.}
Which restrictions on representation and translation complexity preserve rich concepts while preventing the interpreter from supplying the whole solution? A good answer should not privilege one architecture by definition. It should expose a testable relation among description, use, and acquisition cost.

\paragraph{Approximate correction.}
Can evidence transport preserve calibrated uncertainty when filtering is replaced by probabilistic updating or robust model sets? A small numerical error in observations can cause a discontinuous change in an exact candidate set. The metric transfer bound alone does not settle this question.

\paragraph{Acquiring new distinctions.}
Given a failed continuation and an old quotient, what admissible observation supplies the least additional information needed for the enlarged contract? Proposition~\ref{prop:expansion} gives a finite coding bound, not a theorem that the required sensor or experiment exists. Observation cost, disturbance, and identifiability must enter the construction.

\paragraph{Useful explanation across different recipients.}
When do two recipients with different background representations benefit from one artifact, and when is a family of explanations necessary? The counting bound gives only a necessary information condition. A fuller account needs interpreter complexity, partial knowledge, and the recipient's ability to verify and correct the explanation.

These questions define a continuation of the paper rather than a claim that its first draft has completed a universal theory. They focus the next work on places where the existing mathematics identifies a real missing relation.

\section{Conclusion}\label{sec:conclusion}
A correct answer is an achievement. Reusable understanding concerns what remains available afterward: a method for new continuations, a translation another reasoner can use, and a way to recognize when new evidence is needed.

The instrument makes the progression concrete. A display value answers the present query. Two coordinates support every continuation under two operations. A third register becomes indispensable when a new operation is admitted. A translation can preserve the old behavior while a coarsened evidence language prevents a later correction. These are different computational claims.

The proposal is to study representations whose operations, use, transfer, and correction remain explicit. Minimality depends on future questions. Mechanistic attribution depends on permissible evidence. Explanation depends on a recipient and its background. Corrective transfer must preserve the distinctions that can disqualify a model. The encoder, decoder, and translator remain inside the resource account.

Different realizations can meet these obligations without reproducing one human notation. What they cannot do is make fluency, accuracy, formal validity, and structural correspondence interchangeable evidence. A representation becomes intellectually useful when its retained relationships support further questions and its limits remain available to the next act of reasoning.

\clearpage\appendix
\section{A complete trace of the instrument}\label{app:trace}
Start from $(a,b,c)=(1,3,2)$. The first display is four. The base record is $(s,b)=(4,3)$. Apply $H$, then $W$, then the newly admitted $J$.

\begin{center}
\begin{tabular}{@{}lcccl@{}}\toprule
Event & Complete state & Display & Base record & Status\\\midrule
Initial & $(1,3,2)$ & $4$ & $(4,3)$ & $H/W$ contract adequate\\
After $H$ & $(4,3,2)$ & $2$ & $(2,3)$ & Exact base update\\
After $W$ & $(3,4,2)$ & $2$ & $(2,4)$ & Exact base update\\
Before predicting $J$ & $(3,4,2)$ & $2$ & $(2,4)$ & Missing $c$ distinction\\
After $J$ & $(0,4,2)$ & $4$ & $(4,4)$ & New reading acquired\\\bottomrule
\end{tabular}
\end{center}

The last base record does not by itself retain all newly relevant information. The previous display was two and the new display is four, so $c=4-2=2$. The post-operation first register is $4-4=0$. The enlarged display-centered record is therefore $(s,b,c)=(4,4,2)$.

A subsequent $H$ gives $(s,b,c)=(3,4,2)$ and display three. A subsequent $W$ leaves the display at three but changes the second coordinate to $3-4=4$ modulo five. A subsequent $J$ changes the display to zero. Each assertion follows from the enlarged update equations, not from extrapolating a list of displayed numbers.

\subsection{What was actually learned}
The original prediction method already encoded the effect of $H$ and $W$. The $J$ observation supplied a new state distinction. It did not establish the $J$ transition law from scratch, because that law was part of the expanded analytical contract. Learning that law would require a candidate family and experiments capable of distinguishing it from alternatives.

This separation is important. Acquiring a latent value under a known rule, inferring an unknown rule, and transferring a known rule into a new notation are different research achievements. A single successful trace should not be counted as all three.

\subsection{A correction record for the trace}
The record should retain the old scope, the new operation, the failed sufficiency condition, the acquired observation, and the exact post-intervention state reconstruction. It should also preserve the unchanged base equations. A concise corrected explanation would say that the third coordinate is irrelevant to $H/W$ continuations but relevant to any contract admitting $J$.

Nothing in that explanation establishes a hidden psychological state of an agent. It establishes the content and limits of a reusable mathematical model and gives a recipient an explicit way to use and revise it.

\section{Companion verification and draft boundaries}\label{app:verification}
The companion program \code{code/finite\_checks.py} uses integers and exact rational arithmetic. It checks the register construction, partition refinement against a separately coded pair-search procedure, finite transcript discrimination, a stochastic quotient, transport-error recurrences, explanation capacity, informative model filtering, and correction transport. The program enumerates small domains stated separately for each check group. It does not train a model or test an empirical claim about machine cognition.

The finite quotient checks compare the refinement procedure with direct search for a distinguishing word for every small machine in the stated domain. The register checks compare microscopic updates with independently written reduced coordinate updates. The correction checks enumerate maps, candidate sets, and possible evidence, including evidence inconsistent with the current set. Negative controls deliberately use a wrong swap update, a collapsed evidence channel, and a delayed disagreement beyond a fixed test horizon.

The script's result is evidence about the supplied code and finite constructions. The prose arguments supply the general conditional results. A clean compilation, an exact finite check, an independent proof review, and a formal-kernel build remain different kinds of evidence. This draft claims the first two, within the recorded review scope, and does not claim the latter two.

The source package records the actual test output, source hashes, build diagnostics, reference-review scope, and known limitations. Diagrams are generated from mathematical constructions. No synthetic image is offered as an empirical observation. The PDF's version number identifies this first draft, not a completed priority audit or an accepted publication.

\section{Claim map}\label{app:claims}
\begin{center}\small
\begin{tabularx}{\textwidth}{@{}P{1.1in}P{1.55in}X@{}}\toprule
Statement & Main assumption & Conclusion and boundary\\\midrule
Theorem~\ref{thm:quotient} & Finite deterministic operations and all declared continuations. & Minimal behavioral quotient. Does not minimize acquisition or runtime.\\
Proposition~\ref{prop:expansion} & Expanded contract on the same complete state space. & Necessary supplementary capacity. Sharp reuse construction has an explicit refinement condition.\\
Proposition~\ref{prop:identification} & Fixed complete transcript laws and equal priors. & Optimal binary identification bound. New evidence changes the laws.\\
Proposition~\ref{prop:lumpability} & Equal projected rows on each fibre. & Exact stochastic history preservation for reduced-history policies.\\
Theorem~\ref{thm:transport} & Matched operations and query decoders. & Exact forward transfer. Does not grant reverse recovery or target coverage.\\
Proposition~\ref{prop:error} & Uniform residuals and Lipschitz target dynamics. & Finite-horizon error budget for matched operations.\\
Proposition~\ref{prop:explanation} & One fixed recipient background and fixed-length messages. & Necessary communication capacity, not efficient decoding.\\
Proposition~\ref{prop:queries} & Exact realizable candidate family and separating experiments. & Finite informative elimination and decision-tree lower bound.\\
Theorem~\ref{thm:correction} & A commuting evidence square, including contradictory observations. & Filtering inclusion. Uniform equality exactly requires reflection of attainable outcomes.\\\bottomrule
\end{tabularx}\end{center}

All nine entries have proofs in the text. Their mathematical ingredients are classical or elementary. The integrated adequacy record is a design proposal. The proposed account of computational understanding and the open questions are not additional theorems.

\clearpage
\phantomsection\addcontentsline{toc}{section}{References}
\begin{thebibliography}{99}\small
\bibitem{societies} Komorebi P. \emph{Machine Societies: A Theory of Persistent Plural Computational Organization}. Operation Komorebi, version 1.0.0, 2026. Companion theoretical monograph.
\bibitem{jacovi} Alon Jacovi and Yoav Goldberg. Towards Faithfully Interpretable NLP Systems: How Should We Define and Evaluate Faithfulness? \emph{Proceedings of ACL}, pp. 4198--4205, 2020. \url{https://aclanthology.org/2020.acl-main.386/}.
\bibitem{shalizi} Cosma Rohilla Shalizi and James P. Crutchfield. Computational Mechanics: Pattern and Prediction, Structure and Simplicity. \emph{Journal of Statistical Physics}, 104, 2001. \url{https://doi.org/10.1023/A:1010388907793}.
\bibitem{givan} Robert Givan, Thomas Dean, and Matthew Greig. Equivalence Notions and Model Minimization in Markov Decision Processes. \emph{Artificial Intelligence}, 147:163--223, 2003. \url{https://engineering.purdue.edu/~givan/papers/mm.pdf}.
\bibitem{tishby} Naftali Tishby, Fernando C. Pereira, and William Bialek. The Information Bottleneck Method. \emph{Proceedings of the 37th Allerton Conference}, 1999. Author preprint, 2000. \url{https://arxiv.org/abs/physics/0004057}.
\bibitem{beckers} Sander Beckers and Joseph Y. Halpern. Abstracting Causal Models. \emph{Proceedings of AAAI}, 2019. \url{https://arxiv.org/abs/1812.03789}.
\bibitem{geiger} Atticus Geiger, Duligur Ibeling, Amir Zur, Maheep Chaudhary, Sonakshi Chauhan, Jing Huang, Aryaman Arora, Zhengxuan Wu, Noah Goodman, Christopher Potts, and Thomas Icard. Causal Abstraction: A Theoretical Foundation for Mechanistic Interpretability. \emph{Journal of Machine Learning Research}, 26(83):1--64, 2025. \url{https://jmlr.org/papers/v26/23-0058.html}.
\bibitem{hewitt} John Hewitt and Percy Liang. Designing and Interpreting Probes with Control Tasks. \emph{Proceedings of EMNLP-IJCNLP}, pp. 2733--2743, 2019. \url{https://aclanthology.org/D19-1275/}.
\bibitem{sutter} Denis Sutter, Julian Minder, Thomas Hofmann, and Tiago Pimentel. The Non-Linear Representation Dilemma: Is Causal Abstraction Enough for Mechanistic Interpretability? Preprint, 2025. Version 1 consulted. \url{https://arxiv.org/abs/2507.08802v1}.
\bibitem{miller} Tim Miller. Explanation in Artificial Intelligence: Insights from the Social Sciences. \emph{Artificial Intelligence}, 267:1--38, 2019. Author preprint: \url{https://arxiv.org/abs/1706.07269}.
\bibitem{gentner} Dedre Gentner. Structure-Mapping: A Theoretical Framework for Analogy. \emph{Cognitive Science}, 7:155--170, 1983. Author-hosted copy: \url{https://groups.psych.northwestern.edu/gentner/papers/Gentner83.2b.pdf}.
\bibitem{angluin} Dana Angluin. Learning Regular Sets from Queries and Counterexamples. \emph{Information and Computation}, 75:87--106, 1987. \url{https://doi.org/10.1016/0890-5401(87)90052-6}.
\bibitem{blackwell} David Blackwell. Equivalent Comparisons of Experiments. \emph{Annals of Mathematical Statistics}, 24(2):265--272, 1953. \url{https://doi.org/10.1214/aoms/1177729032}.
\bibitem{rischel} Eigil F. Rischel and Sebastian Weichwald. Compositional Abstraction Error and a Category of Causal Models. \emph{Proceedings of UAI}, PMLR 161:1013--1023, 2021. \url{https://proceedings.mlr.press/v161/rischel21a.html}.
\end{thebibliography}
\end{document}
